\documentclass[11pt]{article}
\usepackage[utf8]{inputenc}
\usepackage[T1]{fontenc}
\usepackage[ngerman]{babel}
\usepackage[margin=2.5cm]{geometry}

\title{Algorithm Engineering\\Woche 2, Aufgabe 2}
\author{Jan Vogt}
\date{}

\begin{document}

\maketitle

\section*{Die Abbildung "`Performance gains after Moore's law ends"'}

Die Abbildung auf der ersten Seite des Papers "`There's plenty of room at the Top"' von Leiserson
et al. (Science, 2020) fasst die Aussage des ganzen Papers zusammen. Sie besteht aus zwei
Bereichen. Oben steht ein Kasten "`The Top"' mit den drei Spalten Software, Algorithms und
Hardware architecture. Für jede Spalte gibt es ein Bild (Technology), eine Möglichkeit, wie man
dort noch Leistung herausholen kann (Opportunity), und zwei Beispiele (Examples). Unten steht
"`The Bottom"', gezeichnet als Trichter aus Streifen, die nach unten immer schmaler werden. So wie
ich das verstehe, soll das zeigen, dass dort immer weniger Platz für Verbesserungen übrig ist.

\subsection*{The Bottom}

Mit Bottom ist die Halbleitertechnik gemeint, also das Verkleinern der Transistoren nach Moore's
Law. Über 50 Jahre kam fast die ganze Leistungssteigerung von hier, und zwar automatisch: Ein
Programm wurde auf einem neuen Prozessor schneller, ohne dass man etwas daran ändern musste. Laut
Paper ist dieser Platz aber fast aufgebraucht. Intel hat zum Beispiel die 10-nm-Technik erst 2019
statt 2016 geschafft, und bei ungefähr 5 nm wird die Verkleinerung wohl aufhören, weil sie danach
kaum noch etwas bringt und immer teurer wird. Deshalb sagt die Bildunterschrift, dass sich der
bisherige Trend umdreht und die Gewinne in Zukunft vor allem aus dem Top kommen.

\subsection*{Die drei Spalten im Top}

\textbf{Software.} Die Möglichkeit ist Software performance engineering, also Programme so
umzubauen, dass sie schneller laufen. Das heißt einerseits Software bloat entfernen, weil Software
bisher eher schnell fertig als schnell sein sollte, und andererseits die Software an die Hardware
anpassen, zum Beispiel an mehrere Kerne, Vektoreinheiten und Caches. Im Paper wird eine
$4096 \times 4096$ Matrixmultiplikation in Python, die etwa 7 Stunden braucht, Schritt für Schritt
bis auf 0,41 Sekunden beschleunigt. Das ist mehr als 60.000 mal schneller.

\textbf{Algorithms.} Bessere Algorithmen lösen ein Problem grundsätzlich schneller. Beim
Maximum-Flow-Problem haben neue Algorithmen über 38 Jahre fast so viel gebracht wie Moore's Law,
seit 20 Jahren aber nichts mehr. Fortschritt kommt laut Paper vor allem durch neue Problembereiche
wie Machine Learning (New problem domains) und durch neue Maschinenmodelle, weil die meisten
Algorithmen noch für einen Prozessor entworfen wurden, der nur eine Operation gleichzeitig macht und
bei dem jeder Speicherzugriff gleich lange dauert (New machine models).

\textbf{Hardware architecture.} Wenn es nicht mehr Transistoren gibt, muss man die vorhandenen
besser nutzen (Hardware streamlining). Bei Processor simplification werden aufwendige Kerne durch
einfachere ersetzt, damit mehr Kerne auf den Chip passen. Das lohnt sich aber nur, wenn das
Programm genug Parallelität hat. Bei Domain specialization wird Hardware für einen bestimmten
Bereich gebaut, zum Beispiel GPUs oder die TPU von Google für Deep Learning.

\subsection*{Gesamtaussage}

Auch nach dem Ende von Moore's Law kann man Computer also noch schneller machen, nur nicht mehr
automatisch. Beim Top muss man für jedes Problem gezielt etwas ändern, und die Gewinne kommen laut
Paper ungleichmäßig und werden kleiner, je besser ein Problem schon erforscht ist.

Das ist mir auch im Zusammenhang mit der Vorlesung aufgefallen: Die Spalte Software ist im Grunde
genau das, was wir in Algorithm Engineering machen. Bei der Matrixmultiplikation im Paper kommen
die großen Sprünge durch Parallelisierung, Cache-Ausnutzung und Vektorisierung, also genau die
Themen aus der Kursübersicht. Außerdem steht im Bild für Software die Folge
\texttt{01010011 01100011 01101001 01100101 01101110 01100011 01100101 00000000}. Liest man die
Bytes als ASCII wie in Kapitel 1 von CS:APP, ergibt das das Wort "`Science"' mit einem Null-Byte
am Ende, also genau wie ein String in C.

\end{document}
